\documentclass[11pt,a4paper]{article}

% ---------- Paquets ----------
\usepackage[utf8]{inputenc}
\usepackage[T1]{fontenc}
\usepackage[french]{babel}
\usepackage{lmodern}
\usepackage[margin=2.5cm]{geometry}
\usepackage{xcolor}
\usepackage{booktabs}
\usepackage{tabularx}
\usepackage{enumitem}
\usepackage{fancyhdr}
\usepackage{titlesec}
\usepackage{pgfplots}
\pgfplotsset{compat=1.17}
\usepackage[hidelinks]{hyperref}

% ---------- Couleurs et style ----------
\definecolor{bleu}{HTML}{1F3A5F}
\definecolor{gris}{HTML}{5A6270}

\titleformat{\section}{\Large\bfseries\color{bleu}}{\thesection.}{0.6em}{}
\titleformat{\subsection}{\large\bfseries\color{bleu}}{\thesubsection}{0.6em}{}

\pagestyle{fancy}
\fancyhf{}
\fancyhead[L]{\small\color{gris}Accidents de la route 2023 et 2024}
\fancyhead[R]{\small\color{gris}EPSI, M1}
\fancyfoot[C]{\small\color{gris}\thepage}
\renewcommand{\headrulewidth}{0.4pt}


\frenchsetup{StandardLists=true}
\setlist{itemsep=2pt, topsep=4pt}
\setlist[itemize]{label=\textbullet}

\begin{document}

% ---------- Page de titre ----------
\begin{titlepage}
  \centering
  \vspace*{2cm}
  {\color{gris}\large EPSI, Master 1\par}
  \vspace{1.5cm}
  {\color{bleu}\rule{\linewidth}{1pt}\par}
  \vspace{0.6cm}
  {\Huge\bfseries\color{bleu} Les accidents de la route\\[0.3cm] en France en 2023 et 2024\par}
  \vspace{0.5cm}
  {\Large Rapport sur les données utilisées\par}
  \vspace{0.6cm}
  {\color{bleu}\rule{\linewidth}{1pt}\par}
  \vspace{2cm}

  {\large\bfseries Membres du groupe\par}
  \vspace{0.4cm}
  {\large
    Mohamed Amine BAHROUI\\[2pt]
    Johan BILE\\[2pt]
    Aya HANDIRI\\[2pt]
    Glody KUTUMBAKANA\\[2pt]
    Théo LORIO\\[2pt]
    Raji RAHOUI\par}

  \vfill
  {\color{gris}\today\par}
\end{titlepage}

\tableofcontents
\newpage

% ---------- 1. Introduction ----------
\section{Introduction}

Ce rapport présente les données avec lesquelles nous avons travaillé : des
fichiers publics qui décrivent les accidents de la route en France pour les
années \textbf{2023} et \textbf{2024}.

Notre travail avait trois objectifs simples :
\begin{enumerate}
  \item \textbf{vérifier} que les données sont complètes et cohérentes ;
  \item \textbf{calculer} un indicateur mensuel : le nombre d'accidents par mois ;
  \item \textbf{décider} si ce résultat peut être publié tel quel.
\end{enumerate}

\subsection*{L'essentiel}

\begin{itemize}
    \item Environ \textbf{54\,600 accidents corporels} par an, soit près de
          \textbf{110\,000} sur les deux années.
    \item Les données sont de \textbf{bonne qualité} : pas de doublon, pas de
          date fausse, tous les fichiers sont bien reliés entre eux.
    \item Le nombre d'accidents est \textbf{presque stable} entre 2023 et 2024
          ($-0{,}8\,\%$).
    \item Environ \textbf{3\,200 accidents mortels} par an, soit près de
          6\,\% des accidents.
    \item Décision : \textbf{publication en attente} (\og NO-GO \fg{}) tant que
          nos chiffres n'ont pas été comparés aux chiffres officiels et que la
          complétude des fichiers n'est pas confirmée.
\end{itemize}

% ---------- 2. Les données ----------
\section{D'où viennent les données ?}

Les données viennent des \textbf{fichiers BAAC} (Bulletins d'Analyse des
Accidents Corporels). Chaque fois qu'un accident fait au moins une victime
(blessé ou tué), les forces de l'ordre remplissent une fiche. Ces fiches sont
ensuite publiées en données ouvertes par l'État.

\medskip
\noindent\textbf{Important :} ces fichiers ne contiennent \emph{que} les
accidents avec des victimes. Les accidents avec uniquement des dégâts matériels
n'y sont pas.

\subsection{Quatre fichiers par année}

Pour chaque année, les données sont réparties dans quatre fichiers. On peut les
voir comme quatre « feuilles » qui décrivent le même accident sous des angles
différents.

\begin{center}
\begin{tabularx}{\linewidth}{@{}l X l@{}}
  \toprule
  \textbf{Fichier} & \textbf{Ce qu'il décrit} & \textbf{1 ligne =} \\
  \midrule
  \texttt{caract}    & Les circonstances : date, heure, lumière, météo, lieu (commune, GPS), type de collision & 1 accident \\
  \texttt{lieux}     & La route : type de route, nombre de voies, état de la chaussée, vitesse maximale autorisée & 1 route concernée \\
  \texttt{vehicules} & Les véhicules : catégorie (voiture, moto, vélo\dots), obstacle heurté, point de choc, manœuvre & 1 véhicule \\
  \texttt{usagers}   & Les personnes : place, conducteur ou passager, sexe, année de naissance, équipement de sécurité, \textbf{gravité} & 1 personne \\
  \bottomrule
\end{tabularx}
\end{center}

\subsection{Comment les fichiers sont reliés}

Les quatre fichiers partagent une même colonne : \textbf{\texttt{Num\_Acc}},
le numéro unique de l'accident. C'est grâce à ce numéro qu'on peut retrouver,
pour un accident donné, sa route, ses véhicules et les personnes impliquées.

\subsection{Taille des données}

\begin{center}
\begin{tabular}{@{}l r r r@{}}
  \toprule
  \textbf{Fichier} & \textbf{2023} & \textbf{2024} & \textbf{Colonnes} \\
  \midrule
  \texttt{caract}    &  54\,822 &  54\,402 & 15 \\
  \texttt{lieux}     &  70\,860 &  70\,248 & 18 \\
  \texttt{vehicules} &  93\,585 &  92\,678 & 11 \\
  \texttt{usagers}   & 125\,789 & 125\,187 & 16 \\
  \bottomrule
\end{tabular}
\end{center}

En moyenne, un accident implique environ \textbf{1,7 véhicule} et
\textbf{2,3 personnes}.

\subsection{La gravité : la colonne la plus importante}

Dans le fichier \texttt{usagers}, la colonne \texttt{grav} indique l'état de
chaque personne après l'accident. Attention, les codes ne sont \emph{pas}
rangés du moins grave au plus grave :

\begin{center}
\begin{tabular}{@{}c l r r@{}}
  \toprule
  \textbf{Code} & \textbf{Signification} & \textbf{Personnes 2023} & \textbf{Personnes 2024} \\
  \midrule
  1 & Indemne (pas blessé)  & 53\,399 & 52\,920 \\
  4 & Blessé léger          & 49\,603 & 49\,709 \\
  3 & Blessé hospitalisé    & 19\,271 & 19\,126 \\
  2 & Tué                   &  3\,398 &  3\,432 \\
  $-1$ & Non renseigné      &     118 &       0 \\
  \bottomrule
\end{tabular}
\end{center}

% ---------- 3. Qualité ----------
\section{Qualité des données}

\subsection{Les valeurs manquantes}

Dans ces fichiers, une information absente peut être écrite de plusieurs
façons : case vide, \texttt{-1} ou \texttt{N/A}. Nous avons traité toutes ces
valeurs de la même manière : \emph{non renseignée}.

\paragraph{Fichier \texttt{caract} (accidents).} Il est presque complet. Seule
l'adresse manque parfois (2,5\,\% en 2023 et 4,3\,\% en 2024). Les
coordonnées GPS sont toujours présentes.

\paragraph{Fichier \texttt{lieux} (routes).} Quelques colonnes sont presque
toujours vides, car elles ne servent que dans des cas précis :
\begin{itemize}
  \item \texttt{lartpc}, la largeur du terre-plein central (plus de 99\,\%
        vide), n'existe que sur les routes séparées en deux ;
  \item \texttt{v2}, le complément du numéro de route (environ 92\,\% vide) ;
  \item \texttt{larrout}, la largeur de la route (69 à 74\,\% vide) ;
  \item \texttt{pr} et \texttt{pr1}, les repères kilométriques (39 à 42\,\%
        vides), n'existent que sur les grands axes.
\end{itemize}
La vitesse maximale autorisée (\texttt{vma}) manque dans environ 6\,\% des cas.

\paragraph{Fichier \texttt{vehicules}.} Il est presque complet. Seule la
colonne \texttt{occutc} (nombre de passagers d'un bus ou d'un car) est vide à
99\,\%, ce qui est normal : elle ne concerne que les transports en commun.

\paragraph{Fichier \texttt{usagers} (personnes).} Certaines colonnes sont très
souvent vides, mais c'est \textbf{normal} :
\begin{itemize}
  \item \texttt{etatp}, \texttt{locp}, \texttt{actp} (environ 50 à 92\,\% vides)
        ne concernent que les \textbf{piétons} ;
  \item \texttt{secu2} et \texttt{secu3} (43 à 98\,\% vides) servent à noter un
        2\textsuperscript{e} ou 3\textsuperscript{e} équipement de sécurité,
        ce qui est rare.
\end{itemize}
Les colonnes vraiment utiles (sexe, année de naissance, gravité) manquent dans
moins de 2,1\,\% des cas.

\subsection{Les contrôles réalisés}

\begin{center}
\begin{tabularx}{\linewidth}{@{}X c c@{}}
  \toprule
  \textbf{Contrôle} & \textbf{2023} & \textbf{2024} \\
  \midrule
  Numéros d'accident absents ou en double              & 0 & 0 \\
  Dates impossibles (ex. 31 février) ou de la mauvaise année & 0 & 0 \\
  Lignes de détail sans accident correspondant         & 0 & 0 \\
  Accidents sans aucun détail (route, véhicule, personne) & 0 & 0 \\
  Personnes avec une gravité non renseignée            & 118 & 0 \\
  Mois présents dans le fichier                        & janv. à déc. & janv. à déc. \\
  Accidents écartés du calcul                          & 0 & 0 \\
  \bottomrule
\end{tabularx}
\end{center}

Ces contrôles sont aussi vérifiés \textbf{automatiquement} à la fin du calcul :
chaque accident compté a un numéro unique, une date valide et une gravité
connue. Tous les contrôles sont passés.

\subsection{Le journal des anomalies}

Nous n'avons \textbf{rien inventé} : une gravité inconnue n'a pas été
remplacée par une valeur au hasard. Chaque cas a été noté dans un fichier,
\texttt{journal\_corrections.csv}. Pour chaque cas, on y trouve :
\begin{itemize}
  \item l'année et le numéro de l'accident ;
  \item la ligne exacte dans le fichier d'origine, pour retrouver le cas facilement ;
  \item la règle non respectée, l'action prise et le statut (« à vérifier »).
\end{itemize}

Ce journal contient \textbf{118 lignes} pour 2023, une par personne sans
gravité. Ces 118 personnes appartiennent à \textbf{116 accidents}. Dans chacun
de ces accidents, au moins une autre personne a une gravité connue : aucun
accident n'a donc été perdu pour le calcul. Le journal ne corrige rien dans
les fichiers d'origine : il sert seulement à garder une trace.

% ---------- 4. Résultats ----------
\section{Résultats : le nombre d'accidents par mois}

\begin{center}
\begin{tikzpicture}
\begin{axis}[
  width=\linewidth, height=7cm,
  ybar, bar width=6pt,
  ymin=0, ymax=6000,
  xtick=data,
  symbolic x coords={Jan,Fév,Mar,Avr,Mai,Juin,Juil,Août,Sep,Oct,Nov,Déc},
  ylabel={Nombre d'accidents},
  legend style={at={(0.5,1.03)}, anchor=south, legend columns=2, draw=none, /tikz/every even column/.append style={column sep=12pt}},
  ymajorgrids, grid style={gray!25},
  axis line style={gray!60},
  tick label style={font=\small},
  scaled y ticks=false,
  yticklabel style={/pgf/number format/.cd, 1000 sep={\,}},
]
\addplot[fill=bleu!45, draw=none] coordinates {
  (Jan,4053) (Fév,3682) (Mar,3998) (Avr,4162) (Mai,4767) (Juin,5452)
  (Juil,4754) (Août,4121) (Sep,5161) (Oct,5389) (Nov,4833) (Déc,4450)};
\addplot[fill=bleu, draw=none] coordinates {
  (Jan,4069) (Fév,3792) (Mar,4370) (Avr,4586) (Mai,4517) (Juin,5135)
  (Juil,5209) (Août,4450) (Sep,4620) (Oct,4803) (Nov,4477) (Déc,4374)};
\legend{2023, 2024}
\end{axis}
\end{tikzpicture}
\end{center}

\begin{center}
\begin{tabular}{@{}l r r r r@{}}
  \toprule
  \textbf{Mois} & \textbf{2023} & \textbf{2024} & \textbf{Écart} & \textbf{Écart (\%)} \\
  \midrule
  Janvier   & 4\,053 & 4\,069 &  $+16$ &  $+0{,}4$ \\
  Février   & 3\,682 & 3\,792 & $+110$ &  $+3{,}0$ \\
  Mars      & 3\,998 & 4\,370 & $+372$ &  $+9{,}3$ \\
  Avril     & 4\,162 & 4\,586 & $+424$ & $+10{,}2$ \\
  Mai       & 4\,767 & 4\,517 & $-250$ &  $-5{,}2$ \\
  Juin      & 5\,452 & 5\,135 & $-317$ &  $-5{,}8$ \\
  Juillet   & 4\,754 & 5\,209 & $+455$ &  $+9{,}6$ \\
  Août      & 4\,121 & 4\,450 & $+329$ &  $+8{,}0$ \\
  Septembre & 5\,161 & 4\,620 & $-541$ & $-10{,}5$ \\
  Octobre   & 5\,389 & 4\,803 & $-586$ & $-10{,}9$ \\
  Novembre  & 4\,833 & 4\,477 & $-356$ &  $-7{,}4$ \\
  Décembre  & 4\,450 & 4\,374 &  $-76$ &  $-1{,}7$ \\
  \midrule
  \textbf{Total} & \textbf{54\,822} & \textbf{54\,402} & $\mathbf{-420}$ & $\mathbf{-0{,}8}$ \\
  \bottomrule
\end{tabular}
\end{center}

\paragraph{Ce qu'on observe.}
\begin{itemize}
  \item Le total est presque identique d'une année à l'autre.
  \item Février est toujours le mois le plus calme (c'est aussi le plus court).
  \item Les mois les plus chargés sont juin, septembre et octobre en 2023, et
        juin et juillet en 2024.
  \item D'un mois à l'autre, les écarts entre les deux années peuvent
        atteindre $\pm 10\,\%$, mais ils se compensent sur l'année.
\end{itemize}

\newpage
\subsection{Les accidents mortels}

Un accident est dit \emph{mortel} si au moins une personne y est tuée. On
compte \textbf{3\,203 accidents mortels en 2023} et \textbf{3\,226 en 2024},
soit environ \textbf{6\,\% des accidents}.

\begin{center}
\begin{tabular}{@{}l r r@{}}
  \toprule
  \textbf{Mois} & \textbf{2023} & \textbf{2024} \\
  \midrule
  Janvier   & 232 & 251 \\
  Février   & 219 & 221 \\
  Mars      & 212 & 255 \\
  Avril     & 224 & 244 \\
  Mai       & 250 & 248 \\
  Juin      & 287 & 283 \\
  Juillet   & 304 & 280 \\
  Août      & 287 & 304 \\
  Septembre & 299 & 258 \\
  Octobre   & 311 & 316 \\
  Novembre  & 287 & 275 \\
  Décembre  & 291 & 291 \\
  \midrule
  \textbf{Total} & \textbf{3\,203} & \textbf{3\,226} \\
  \bottomrule
\end{tabular}
\end{center}

\subsection{Point d'attention : le calcul de la gravité}

Les codes de gravité ne sont pas rangés dans l'ordre (le code 1 veut dire
\og indemne \fg{} et le code 2 \og tué \fg{}). On ne peut donc pas les comparer
comme des nombres : si on gardait le plus petit code, un accident avec une
personne indemne et une personne tuée serait classé \og indemne \fg{}.

Pour éviter cette erreur, nous regardons simplement si \textbf{au moins une
personne a le code 2} (tué). Un contrôle automatique vérifie que ce résultat
correspond bien aux lignes du fichier \texttt{usagers}.

% ---------- 5. Décision ----------
\section{Décision et limites}

\subsection{Peut-on publier ?}

\textbf{Pas encore : NO-GO.} Les données sont propres et tous nos contrôles
sont passés, mais deux éléments nous manquent :
\begin{enumerate}
  \item \textbf{les chiffres officiels publiés}, mois par mois, pour vérifier
        que nous trouvons les mêmes résultats ;
  \item \textbf{une confirmation que les fichiers sont complets}, venant du
        producteur des données. Avoir des accidents de janvier à décembre ne
        prouve pas, à lui seul, qu'il n'en manque aucun.
\end{enumerate}
Le carnet est prêt : il suffit d'y entrer les chiffres officiels et de
confirmer la complétude pour obtenir un \og GO \fg{}. Les résultats pourront
alors être publiés avec le journal des anomalies.

\subsection{Les limites}

\begin{itemize}
  \item \textbf{Nombre d'accidents $\neq$ risque.} Pour mesurer le danger réel,
        il faudrait connaître le trafic (nombre de kilomètres parcourus). Plus
        de voitures sur la route donne plus d'accidents, sans que la route soit
        plus dangereuse.
  \item \textbf{Seulement deux années.} C'est trop peu pour parler de tendance.
  \item \textbf{Accidents matériels absents.} Ils ne sont pas dans les fichiers
        BAAC.
  \item \textbf{Saisie humaine.} Les fiches sont remplies sur le terrain ;
        certaines informations peuvent être approximatives.
\end{itemize}

% ---------- 6. Conclusion ----------
\section{Conclusion}

Les données BAAC 2023 et 2024 sont fiables et bien organisées : quatre fichiers
reliés par un numéro d'accident, sans doublon ni date fausse. Le nombre
d'accidents corporels est resté stable, autour de 54\,600 par an.

Le calcul est en place et vérifié : environ 3\,200 accidents mortels par an.
Il reste à comparer nos résultats aux chiffres officiels et à confirmer que
les fichiers sont complets. Une fois ces deux points réglés, l'indicateur
mensuel pourra être publié en toute confiance.

\end{document}
